% =========================================================
% CHAPTER 3: SYSTEM DESIGN
% =========================================================

\chapter{SYSTEM DESIGN}
\label{chap:system_design}

The design of QueryPilot is organized around a sequential pipeline that
processes each incoming SQL query through a series of well-defined stages,
from initial receipt to final execution and feedback collection. This chapter
describes the flow of control and data through the system using three
complementary flow diagrams: the overall QueryPilot execution flow, the SAGE
decision flow, and the workload drift detection flow.

\section{FLOW DIAGRAMS}
\label{sec:flow_diagrams}

The system design is best understood through three flows, each representing a
distinct operational dimension of QueryPilot. The execution flow describes the
end-to-end path of a single query. The SAGE decision flow describes the logic
by which the decision controller selects an action. The drift detection flow
describes how the Page-Hinkley detector is updated and how a detected shift
triggers a response.

\subsection{QueryPilot Execution Flow}
\label{subsec:execution_flow}

The QueryPilot execution flow describes the complete sequence of operations
performed for every incoming SQL query. The flow is sequential and consists of
ten stages. In the first stage, the Workload Collector receives the SQL query
from the application and records its arrival metadata. In the second stage,
the Feature Extractor runs MySQL's EXPLAIN command on the query and extracts
eight structural features from the resulting execution plan. In the third
stage, the Candidate Generator injects supported optimizer hints into the
original query to produce up to four candidate variants. In the fourth stage,
the Performance Prediction Model predicts the expected execution time for each
candidate using the trained ensemble. In the fifth stage, the Uncertainty
Estimator computes the relative ensemble disagreement for each prediction. In
the sixth stage, the Drift Detector checks whether the global Page-Hinkley
statistic signals a workload shift. In the seventh stage, SAGE applies its
decision logic to select the best candidate or fall back to the native query.
In the eighth stage, the Execution Layer issues the selected SQL to MySQL and
measures client-side latency. In the ninth stage, the Feedback Collector
records the actual execution time and appends it to the training buffer. In the
tenth and final stage, the model retrains if the buffer has reached the minimum
evidence threshold, completing one full iteration of the adaptive loop.

As shown in \figref{fig:execution_flow}, the ten stages form a closed loop
where each query's outcome feeds back into the learning and monitoring
components for subsequent queries.

\begin{figure}[htbp]
  \centering
  \fbox{\parbox{0.85\textwidth}{%
    \centering
    \vspace{1.5cm}
    \textbf{Figure Placeholder}\\[0.3cm]
    QueryPilot Execution Flow\\
    (Replace with actual flow diagram image)
    \vspace{1.5cm}
  }}
  \caption{QueryPilot Execution Flow: Ten-Stage Pipeline}
  \label{fig:execution_flow}
\end{figure}

The flow diagram in \figref{fig:execution_flow} illustrates how a query enters
the system, traverses the pipeline, and returns results to the calling
application while simultaneously updating the learning and monitoring state.

\subsection{SAGE Decision Flow}
\label{subsec:sage_flow}

The SAGE decision flow describes the logic executed by the Stability-Aware
Guarded Execution controller for every incoming query after feature extraction
and candidate generation. SAGE applies three sequential checks. In the first
check, SAGE evaluates whether the performance model is trained and whether the
system has seen at least $n_{min} = 10$ queries. If either condition is not
met, SAGE selects the \textit{gather evidence} action: the native query is
executed and its result is recorded. No optimization is attempted. In the
second check, SAGE evaluates whether the drift detector has flagged an active
workload shift. If drift is detected, SAGE selects the \textit{fallback}
action: the native query is executed, and the training buffer is cleared to
prevent the stale model from being used for future decisions. In the third
check, which applies only when the model is trained and no drift is active,
SAGE evaluates each candidate. Candidates whose relative uncertainty
$u(x) = \sigma(x)/\mu(x)$ exceeds $\tau_u = 0.8$ are excluded. Candidates
whose predicted mean exceeds the native predicted mean by more than
$\tau_r = 0.15$ (15\%) are also excluded. If a non-native candidate survives
both filters and is predicted to be faster than native, SAGE selects the
\textit{adapt} action. Otherwise, SAGE selects \textit{fallback}.

The decision flow is shown in \figref{fig:sage_flow}.

\begin{figure}[htbp]
  \centering
  \fbox{\parbox{0.85\textwidth}{%
    \centering
    \vspace{1.5cm}
    \textbf{Figure Placeholder}\\[0.3cm]
    SAGE Decision Flow\\
    (Replace with actual flowchart image)
    \vspace{1.5cm}
  }}
  \caption{SAGE Decision Flow: Adapt, Gather Evidence, and Fallback Paths}
  \label{fig:sage_flow}
\end{figure}

As shown in \figref{fig:sage_flow}, the three-branch structure ensures that
SAGE never takes an action it cannot justify with sufficient evidence and
acceptable uncertainty.

\subsection{Workload Drift Detection Flow}
\label{subsec:drift_flow}

The workload drift detection flow describes how the Page-Hinkley detector is
updated and how a detected shift triggers SAGE's fallback response. After every
query execution, the actual execution time in milliseconds is passed to the
WorkloadDriftDetector. The detector maintains two sets of Page-Hinkley
statistics: one global detector across all queries, and one per-query detector
keyed by the first 60 characters of the query signature. The global detector
is updated with the new execution time. The Page-Hinkley statistic
$PH_n = \sum_{t=1}^{n}(x_t - \bar{x}_t - \delta) - \min_{1 \le t \le n}
\sum$ is incremented. If $PH_n$ exceeds the threshold $\lambda = 25$, the
detector returns True, signalling drift. On the next query execution, SAGE
reads this drift flag. If the flag is True, SAGE clears the training buffer,
resets to gather-evidence mode, and clears the drift flag. The detector
continues to accumulate statistics on subsequent queries.

The drift detection flow is illustrated in \figref{fig:drift_flow}.

\begin{figure}[htbp]
  \centering
  \fbox{\parbox{0.85\textwidth}{%
    \centering
    \vspace{1.5cm}
    \textbf{Figure Placeholder}\\[0.3cm]
    Workload Drift Detection Flow\\
    (Replace with actual flowchart image)
    \vspace{1.5cm}
  }}
  \caption{Workload Drift Detection Flow: Page-Hinkley Update and Trigger}
  \label{fig:drift_flow}
\end{figure}

As illustrated in \figref{fig:drift_flow}, drift detection operates as a
continuous background process that gates SAGE's adaptation decisions.
